\section{Aim and Research Questions}

% Formulate in one or two sentences what the aim is. Remember to try to push the horizon.

The primary aim of this project is to investigate how Active Learning in primary mathematics education (ages 10--12) can be systematically supported through a responsible, interactive AI system. Specifically, the project examines how a Retrieval-Augmented Generation (RAG) architecture can be designed to function as a \enquote{More Knowledgeable Peer}~\cite{vygotsky1978mind}, providing adaptive scaffolding that aligns with diverse learner profiles. Unlike standard chatbots, a RAG architecture ensures transparency and factual consistency by grounding AI responses in verified mathematical curricula, thereby mitigating the black-box nature of LLM reasoning~\cite{yu2024rag}. By combining adaptive scaffolding driven by a dynamic user model — one that approximates each learner's cognitive profile, including neurodivergent characteristics~\cite{esteban2017how} - with game-based interaction elements to mitigate mathematical anxiety~\cite{lim2012creating}, the system seeks to transform mathematics learning into an engaging, low-stakes, and cognitively supportive experience.

% Formulate 3-4 specific research questions that you aim to focus on in the project. The research questions should be concrete and challenging, not too broad, and possible to answer with the methods that you plan to use.

To achieve this aim, the project investigates the following research questions:

\begin{description}

    \item [RQ1:] How can a dynamic user model be operationalized within a RAG framework to approximate a student's Zone of Proximal Development and deliver scaffolding that promotes active reasoning rather than answer dependency~\cite{nirenburg2017cognitive}~\cite{ricci2022go}?

    \item [RQ2:] To what extent can game-based design reframe mathematical errors as constructive learning moments, thereby reducing anxiety and sustaining motivation~\cite{ryan2000self}?

    \item [RQ3:] How can the \enquote{MiniGuide}~\cite{lithner2025miniguide} pedagogical strategy enable the system's generation logic to ensure that AI responses provide transparent guidance rather than direct answers?

\end{description}